\documentclass[12pt, a4paper]{article}
\usepackage[margin=1in]{geometry}
\usepackage{xcolor}
\usepackage{minted}
\usepackage{tcolorbox}
\usepackage{graphicx}
\usepackage{hyperref}
\usepackage{fancyhdr}
\usepackage{caption}
\usepackage{tikz}
\usepackage[utf8]{inputenc}
\usepackage{newunicodechar}

% Fix Unicode character issues
\newunicodechar{⚠}{*}
\newunicodechar{️}{}
\newunicodechar{✓}{[OK]}
\newunicodechar{✗}{[X]}
\newunicodechar{✅}{[OK]}
\newunicodechar{❌}{[ERROR]}

% Define custom colors
\definecolor{darkbg}{HTML}{0a0e27}
\definecolor{codebg}{HTML}{1a1f3a}
\definecolor{accentpurple}{HTML}{7c3aed}
\definecolor{success}{HTML}{10b981}
\definecolor{danger}{HTML}{ef4444}
\definecolor{notebg}{HTML}{f3f4f6}

% Minted configuration for prettier code blocks
\usemintedstyle{native}
\setminted{
    bgcolor=codebg,
    fontsize=\small,
    linenos=false,
    breaklines=true
}

% Custom tcolorbox styles
\newtcolorbox{devnote}[1][]{
    colback=notebg,
    colframe=accentpurple,
    left=10pt,
    right=10pt,
    top=10pt,
    bottom=10pt,
    arc=3pt,
    title=\textbf{Developer's Note},
    title filled=true,
    colbacktitle=accentpurple,
    coltitle=white,
    #1
}

\newtcolorbox{techdetail}[1][]{
    colback=white,
    colframe=danger,
    left=10pt,
    right=10pt,
    top=10pt,
    bottom=10pt,
    arc=3pt,
    title=\textbf{WATCH OUT},
    title filled=true,
    colbacktitle=danger,
    coltitle=white,
    #1
}

\newtcolorbox{highlight}[1][]{
    colback=white,
    colframe=success,
    left=10pt,
    right=10pt,
    top=10pt,
    bottom=10pt,
    arc=3pt,
    #1
}

% Header and Footer
\pagestyle{fancy}
\fancyhf{}
\setlength{\headheight}{14pt}
\rhead{\textit{AI Spam Guardian}}
\lhead{\textit{Technical Deep Dive}}
\cfoot{\thepage}

\title{\textbf{Building a Modern Spam Detection System with Flask and Scikit-Learn}}
\author{\textbf{Vishawakarma University} \\[0.5em] {\large \textit{Prn: 31232431}} \\[0.3em] Developer's Technical Blog}
\date{\today}

\begin{document}

\maketitle

\begin{abstract}
Over the past couple of sprints, I've been working on an interesting real-world ML problem: building a production-ready spam detection system with a clean API and snappy frontend. This post walks through the entire stack---from model training to React-style state management in vanilla JavaScript. I'll share some decisions I made, pitfalls I hit, and why I chose certain technologies over others.
\end{abstract}

\tableofcontents
\newpage

\section{The Big Picture}

Let me start by painting the overall architecture. We're building a \textbf{three-layer stack}:

\begin{enumerate}
    \item \textbf{Machine Learning Layer}: A Naive Bayes classifier trained on the UCI spam dataset (5,572 messages total)
    \item \textbf{API Layer}: Flask backend serving predictions via REST endpoints
    \item \textbf{UI Layer}: A modern, dark-mode dashboard with real-time charts and instant feedback
\end{enumerate}

The whole thing runs locally on a single development machine, but the architecture is clean enough that scaling to a production environment would mostly involve containerization and maybe switching from Flask's development server to Gunicorn.

\subsection{Why This Stack?}

Here's the honest truth: \textbf{scikit-learn is still one of the best libraries} for getting a classifier into production quickly. Yes, transformer models are sexy, but for a binary classification task with a modest dataset, Multinomial Naive Bayes gives us 98\% accuracy without breaking a sweat. I could've used a deep learning framework, but that felt like overkill---and debugging neural networks in production is nobody's idea of fun.

\section{Methodology and System Architecture}

Our approach follows a classic machine learning pipeline with three distinct phases: \textbf{data preparation}, \textbf{model training}, and \textbf{inference serving}. Let me break down the methodology:

\subsection{Process Flow Diagram}

Here's how data flows through the system:

\begin{center}
\begin{tikzpicture}[
    node distance=2.5cm,
    box/.style={rectangle, draw=accentpurple, line width=2pt, minimum width=3cm, minimum height=0.8cm, text centered, font=\textbf},
    process/.style={rectangle, draw=success, line width=1.5pt, minimum width=3.5cm, minimum height=0.8cm, text centered},
    io/.style={rectangle, draw=danger, line width=1.5pt, minimum width=3cm, minimum height=0.8cm, text centered},
    arrow/.style={->, line width=2pt, color=accentpurple}
]

    \node[io] (input) {Input Message};
    \node[box, below=1.5cm of input] (preprocess) {Text Preprocessing};
    \node[process, below=1.5cm of preprocess] (vectorize) {CountVectorizer};
    \node[box, below=1.5cm of vectorize] (predict) {Naive Bayes Model};
    \node[io, below=1.5cm of predict] (output) {Spam/Ham + Confidence};
    
    \draw[arrow] (input) -- (preprocess);
    \draw[arrow] (preprocess) -- (vectorize);
    \draw[arrow] (vectorize) -- (predict);
    \draw[arrow] (predict) -- (output);
    
    % Training phase (right side)
    \node[io, right=5cm of input] (dataset) {UCI Spam Dataset};
    \node[box, below=1.5cm of dataset] (train_split) {80/20 Train/Test Split};
    \node[process, below=1.5cm of train_split] (train_vec) {Fit CountVectorizer};
    \node[box, below=1.5cm of train_vec] (train_model) {Train Naive Bayes};
    \node[process, below=1.5cm of train_model] (eval) {Evaluate (98.39\%)};
    
    \draw[arrow] (dataset) -- (train_split);
    \draw[arrow] (train_split) -- (train_vec);
    \draw[arrow] (train_vec) -- (train_model);
    \draw[arrow] (train_model) -- (eval);
    
\end{tikzpicture}
\end{center}

\begin{devnote}
    The diagram shows the dual nature: on the left, inference-time predictions; on the right, the training pipeline. In production, these run at different times---training happens once, inference happens constantly.
\end{devnote}

\subsection{Step-by-Step Methodology}

\textbf{Phase 1: Data Acquisition and Cleaning}

\begin{itemize}
    \item Source: UCI Machine Learning Repository (spam.csv)
    \item Format: CSV with two columns (label, message)
    \item Cleaning: Remove null values, standardize encoding to UTF-8
    \item Labels: Binary classification (Ham = 0, Spam = 1)
\end{itemize}

\textbf{Phase 2: Feature Engineering}

\begin{itemize}
    \item Technique: Bag-of-Words (BoW) with CountVectorizer
    \item Tokenization: Split messages into individual words
    \item Vocabulary: Automatically built from training set
    \item Vectorization: Convert text into sparse feature matrices
\end{itemize}

\begin{devnote}
    I initially considered TF-IDF, but CountVectorizer with default parameters actually performed better on this dataset. Sometimes simpler is better.
\end{devnote}

\textbf{Phase 3: Model Selection and Training}

\begin{itemize}
    \item Algorithm: Multinomial Naive Bayes (fast, interpretable, works well with BoW)
    \item Why Naive Bayes? Assumes word independence, which isn't perfectly true but gives excellent results for this domain
    \item Training: On 80\% subset (4,457 messages)
    \item Testing: On 20\% holdout (1,115 messages)
\end{itemize}

\textbf{Phase 4: Inference and Confidence Scoring}

\begin{itemize}
    \item Input: Raw user message
    \item Vectorization: Transform using learned vocabulary
    \item Prediction: Pass through trained model
    \item Confidence: Extract probability from predict\_proba()
    \item Output: Class label + confidence percentage
\end{itemize}

\section{The Machine Learning Foundation}

Let's start with the heart of the system: \texttt{model.py}. This file encapsulates all the ML logic into a clean, object-oriented interface. Here's why I structured it this way...

\subsection{Class Architecture}

\begin{minted}{python}
class SpamModel:
    def __init__(self, data_path=None):
        # Try multiple possible paths for the dataset
        if data_path is None:
            possible_paths = [
                r'C:\Users\Fofan\Downloads\archive\spam.csv',
                r'C:\Users\Fofan\OneDrive\Desktop\NLP_P1\spam.csv',
                'spam.csv',
                os.path.join(os.path.dirname(__file__), 'spam.csv'),
            ]
            
            self.data_path = None
            for path in possible_paths:
                if os.path.exists(path):
                    self.data_path = path
                    break
            
            if self.data_path is None:
                self.data_path = possible_paths[0]
        else:
            self.data_path = data_path
            
        self.model = None
        self.cv = None
        self.metrics = {}
        self.stats = {}
        self.is_trained = False
        self.prediction_counts = {'ham': 0, 'spam': 0}
\end{minted}

\begin{devnote}
I'm checking multiple paths because... let's be real, people store datasets everywhere. By searching through common locations before failing, we reduce setup friction. The tradeoff? A bit of defensive code in the constructor. Worth it.
\end{devnote}

\subsection{Training Pipeline}

The training logic handles all the typical NLP preprocessing steps. Here's what's happening step-by-step:

\begin{minted}{python}
def train(self):
    if not os.path.exists(self.data_path):
        raise FileNotFoundError(...)
    
    # Load and clean
    df = pd.read_csv(self.data_path, encoding='latin-1')
    df = df[['v1', 'v2']]
    df.columns = ['label', 'message']
    
    # Gather stats for the dashboard
    ham_count = int(len(df[df['label'] == 'ham']))
    spam_count = int(len(df[df['label'] == 'spam']))
    
    df['length'] = df['message'].apply(len)
    avg_ham_len = float(df[df['label'] == 'ham']['length'].mean() or 0)
    avg_spam_len = float(df[df['label'] == 'spam']['length'].mean() or 0)
\end{minted}

\subsubsection{Feature Extraction and Model Training}

\begin{minted}{python}
    # Convert labels to binary
    df['label'] = df['label'].map({'ham': 0, 'spam': 1})
    
    # Train/test split
    X_train, X_test, y_train, y_test = train_test_split(
        df['message'], df['label'], test_size=0.2, random_state=42
    )
    
    # Vectorize text into feature space
    self.cv = CountVectorizer()
    X_train_numeric = self.cv.fit_transform(X_train)
    X_test_numeric = self.cv.transform(X_test)
    
    # Train the classifier
    self.model = MultinomialNB()
    self.model.fit(X_train_numeric, y_train)
    
    # Evaluate
    predictions = self.model.predict(X_test_numeric)
    acc = accuracy_score(y_test, predictions)
\end{minted}

\begin{highlight}
    This pipeline is \textbf{synchronous and blocking}---the entire model trains on startup. In production, I'd push this to a background task (Celery, for instance), but for a proof-of-concept, it's perfectly fine.
\end{highlight}

\subsubsection{Storing Metrics}

Here's where I make sure the dashboard gets real numbers:

\begin{minted}{python}
    self.metrics['accuracy'] = float(acc)
    self.stats = {
        'distribution': {'Ham': ham_count, 'Spam': spam_count},
        'avg_length': {'Ham': round(avg_ham_len, 1), 
                       'Spam': round(avg_spam_len, 1)},
        'accuracy': round(float(acc) * 100, 2),
        'total_messages': ham_count + spam_count
    }
    
    self.is_trained = True
    print("Model trained successfully.")
\end{minted}

\begin{techdetail}
    I'm storing accuracy \textbf{twice}---once as a decimal (for internal logic) and once as a percentage string (for the UI). Feels redundant, but it makes the downstream code cleaner.
\end{techdetail}

\subsection{Making Predictions}

\begin{minted}{python}
def predict(self, message):
    if not self.is_trained:
        raise Exception("Model is not trained yet.")
    
    data = self.cv.transform([message])
    result = self.model.predict(data)
    probability = self.model.predict_proba(data)[0].max()
    
    label = "SPAM" if result[0] == 1 else "HAM"
    
    # Track predictions for live stats
    if result[0] == 1:
        self.prediction_counts['spam'] += 1
    else:
        self.prediction_counts['ham'] += 1
    
    return {
        "prediction": label,
        "confidence": round(float(probability) * 100, 2)
    }
\end{minted}

\begin{devnote}
Notice how I'm tracking prediction counts? That's optional instrumentation---later, we use this to show users how many messages the model has classified since the server started. It's a tiny touch, but it makes the interface feel more alive.
\end{devnote}

\section{The Flask API Layer}

Now for the glue: \texttt{app.py}. This is where we expose the model to the web.

\subsection{Initialization and Error Handling}

\begin{minted}{python}
from flask import Flask, request, jsonify, render_template
from model import SpamModel
import os

app = Flask(__name__)
spam_detector = SpamModel()

startup_error = None
try:
    spam_detector.train()
    print("Model trained successfully!")
except Exception as e:
    startup_error = str(e)
    print("Failed to initialize model on startup.")

@app.route('/')
def index():
    return render_template('index.html', error=startup_error)
\end{minted}

Here's my philosophy: \textbf{fail gracefully}. If the dataset is missing, we still serve the UI but show an error message explaining exactly what went wrong. This is way better than a hard crash.

\subsection{The Predict Endpoint}

This is the workhorse of the API:

\begin{minted}{python}
@app.route('/predict', methods=['POST'])
def predict():
    if not spam_detector.is_trained:
        return jsonify({
            "error": "Model failed to train. Check spam.csv",
            "status": "error"
        }), 500
    
    data = request.get_json()
    
    if not data:
        return jsonify({
            "error": "No data provided",
            "status": "error"
        }), 400
    
    message = data.get('message', '').strip()
    
    if not message:
        return jsonify({
            "error": "Message cannot be empty.",
            "status": "error"
        }), 400
    
    if len(message) < 5:
        return jsonify({
            "error": "Message too short",
            "status": "error"
        }), 400
    
    try:
        result = spam_detector.predict(message)
        result['status'] = 'success'
        return jsonify(result)
    except Exception as e:
        print("Error during prediction:")
        return jsonify({
            "error": "Error analyzing message",
            "status": "error"
        }), 500
\end{minted}

\begin{techdetail}
    I'm adding an extra validation layer here even though we validate on the frontend. \textbf{Never trust the client}. Someone could bypass the UI and send raw requests. You need server-side validation.
\end{techdetail}

\subsection{Statistics Endpoint}

\begin{minted}{python}
@app.route('/stats', methods=['GET'])
def stats():
    if not spam_detector.is_trained:
        return jsonify({
            "error": "Model not trained",
            "status": "error"
        }), 500
    
    stats = spam_detector.get_stats()
    stats['status'] = 'success'
    return jsonify(stats)

if __name__ == '__main__':
    app.run(debug=True, threaded=True, 
            host='127.0.0.1', port=5000)
\end{minted}

\begin{devnote}
I'm using \texttt{threaded=True} because we might have concurrent requests. The default single-threaded Flask mode would queue them up, which feels sluggish. For production, this should run under Gunicorn or uWSGI anyway.
\end{devnote}

\section{The Frontend: Vanilla JavaScript at Scale}

This is where things get interesting. I wanted to prove that \textbf{you don't always need React} to build a responsive, state-driven UI. Here's how \texttt{app.js} orchestrates everything:

\subsection{Initialization and DOM Management}

\begin{minted}{javascript}
document.addEventListener('DOMContentLoaded', () => {
    console.log('Spam Guardian Initialized');
    
    // Cache DOM elements
    const form = document.getElementById('predict-form');
    const input = document.getElementById('message-input');
    const submitBtn = document.getElementById('submit-btn');
    const resultContainer = document.getElementById('result-container');
    const charCounter = document.getElementById('char-count');
    
    // Result display elements
    const resultLabel = document.getElementById('result-label');
    const resultDescription = document.getElementById('result-description');
    const resultIcon = document.getElementById('result-icon');
    const confidencePercentage = document.getElementById('confidence-percentage');
    const confidenceFill = document.getElementById('confidence-fill');
    
    // Stats elements
    const hamLenStat = document.getElementById('stat-ham-len');
    const spamLenStat = document.getElementById('stat-spam-len');
    const accuracyStat = document.getElementById('accuracy-stat');
    const totalMessagesStat = document.getElementById('total-messages-stat');
    
    let distributionChart = null;
    let isLoading = false;
\end{minted}

\begin{devnote}
I'm caching DOM references at the top because querying \texttt{getElementById()} repeatedly is wasteful. This is JavaScript 101, but it's \textbf{so easy to forget} when you're rapid-iterating.
\end{devnote}

\subsection{Form Submission Logic}

\begin{minted}{javascript}
if (form) {
    form.addEventListener('submit', async (e) => {
        e.preventDefault();
        
        const message = input.value.trim();
        
        if (!message) {
            showError('Please enter a message to analyze');
            return;
        }
        
        if (message.length < 5) {
            showError('Message must be 5+ characters');
            return;
        }
        
        await analyzeMessage(message);
    });
}
\end{minted}

\subsubsection{The Analyze Function}

\begin{minted}{javascript}
async function analyzeMessage(message) {
    if (isLoading) return; // Prevent double-submission
    
    isLoading = true;
    submitBtn.disabled = true;
    submitBtn.innerHTML = '<span>Analyzing...</span>';
    
    if (window.lucide) lucide.createIcons();
    resultContainer.classList.add('hidden');
    
    try {
        console.log('Sending to /predict...');
        
        const response = await fetch('/predict', {
            method: 'POST',
            headers: { 'Content-Type': 'application/json' },
            body: JSON.stringify({ message })
        });
        
        if (!response.ok) {
            const errorData = await response.json();
            throw new Error(errorData.error || 'Analysis failed');
        }
        
        const data = await response.json();
        console.log('Result:', data);
        
        if (data.prediction && data.confidence !== undefined) {
            displayResult(data.prediction, data.confidence);
        } else {
            throw new Error('Invalid response from server');
        }
        
    } catch (error) {
        console.error('Error:', error);
        showError(error.message || 'Failed to analyze');
    } finally {
        isLoading = false;
        submitBtn.disabled = false;
        submitBtn.innerHTML = '<span>Scan Message</span>';
        if (window.lucide) lucide.createIcons();
    }
}
\end{minted}

\begin{highlight}
    See that \texttt{isLoading} guard at the top? That prevents the user from spamming the button and overwhelming the server. Small detail, massive UX improvement.
\end{highlight}

\subsection{Displaying Results}

\begin{minted}{javascript}
function displayResult(prediction, confidence) {
    console.log('Result: ' + prediction + ' (' 
                + confidence + '%%)');
    
    resultContainer.classList.remove('hidden', 
        'safe-result', 'spam-result');
    
    const isSafe = prediction === 'HAM';
    const resultClass = isSafe ? 'safe-result' 
                                : 'spam-result';
    resultContainer.classList.add(resultClass);
    
    if (isSafe) {
        resultIcon.setAttribute('data-lucide', 
            'check-circle');
        resultLabel.textContent = '[OK] Safe Message';
        resultDescription.textContent = 
            'This message is legitimate and secure.';
    } else {
        resultIcon.setAttribute('data-lucide', 
            'alert-triangle');
        resultLabel.textContent = '[WARN] Spam Detected';
        resultDescription.textContent = 
            'This message is flagged as potential spam.';
    }
    
    confidencePercentage.textContent = confidence + '%';
    confidenceFill.style.width = confidence + '%';
    
    if (window.lucide) lucide.createIcons();
    
    setTimeout(() => {
        resultContainer.scrollIntoView({ 
            behavior: 'smooth', 
            block: 'nearest' 
        });
    }, 100);
}
\end{minted}

\subsection{Fetching and Displaying Statistics}

\begin{minted}{javascript}
async function fetchStats() {
    try {
        console.log('Fetching statistics...');
        
        const response = await fetch('/stats');
        
        if (!response.ok) {
            throw new Error('Failed to fetch stats');
        }
        
        const data = await response.json();
        console.log('Statistics loaded:', data);
        
        // Update stats displays
        if (hamLenStat && data.avg_length) {
            hamLenStat.textContent = 
                data.avg_length.Ham + ' chars';
        }
        if (spamLenStat && data.avg_length) {
            spamLenStat.textContent = 
                data.avg_length.Spam + ' chars';
        }
        
        const totalHam = data.distribution?.Ham || 0;
        const totalSpam = data.distribution?.Spam || 0;
        const totalMessages = totalHam + totalSpam;
        
        // Display real metrics
        if (totalMessagesStat) {
            totalMessagesStat.textContent = 
                totalMessages.toLocaleString();
        }
        if (accuracyStat && data.accuracy !== undefined) {
            accuracyStat.textContent = 
                data.accuracy + '%';
        }
        
        renderChart(totalHam, totalSpam);
        
    } catch (error) {
        console.error('Failed to fetch stats:', error);
    }
}
\end{minted}

\begin{devnote}
That bit at the end---where we display real accuracy instead of placeholder dashes---is the key insight. \textbf{Always surface actual metrics}. Users trust systems that show them verifiable numbers.
\end{devnote}

\subsection{Chart Rendering}

\begin{minted}{javascript}
function renderChart(hamCount, spamCount) {
    try {
        const ctx = document.getElementById('distributionChart');
        
        if (!ctx) {
            console.error('Chart canvas not found');
            return;
        }
        
        console.log('Chart: Ham=' + hamCount + 
                    ', Spam=' + spamCount);
        
        if (distributionChart) {
            distributionChart.destroy();
        }
        
        if (window.Chart) {
            Chart.defaults.font.family = 
                "'Inter', sans-serif";
            Chart.defaults.color = '#a0aec0';
            
            distributionChart = new Chart(ctx, {
                type: 'doughnut',
                data: {
                    labels: ['Safe (Ham)', 'Spam'],
                    datasets: [{
                        data: [hamCount, spamCount],
                        backgroundColor: [
                            'rgba(16, 185, 129, 0.8)',
                            'rgba(239, 68, 68, 0.8)'
                        ],
                        borderColor: [
                            'rgb(16, 185, 129)',
                            'rgb(239, 68, 68)'
                        ],
                        borderWidth: 2
                    }]
                },
                options: {
                    responsive: true,
                    maintainAspectRatio: true,
                    plugins: {
                        legend: {
                            labels: {
                                color: '#f0f4f8',
                                font: { size: 12 }
                            }
                        }
                    }
                }
            });
        }
    } catch (error) {
        console.error('Error rendering chart:', error);
    }
}
\end{minted}

\section{Design Patterns and Architecture Decisions}

\subsection{Why Class-Based ML Wrapper?}

Looking back at the \texttt{SpamModel} class, I wanted to answer a simple question: \textbf{where does state live?} By wrapping the scikit-learn objects in a class, I get:

\begin{itemize}
    \item \textbf{Lazy Initialization}: The model doesn't exist until we call \texttt{train()}
    \item \textbf{Testability}: Easy to instantiate fresh models in unit tests
    \item \textbf{Extensibility}: Later, I can swap in different algorithms without changing the API
\end{itemize}

Plus, it forces me to think about the \textbf{interface} with Flask, not just the implementation.

\subsection{Frontend State Management Without React}

I could've used Redux or Zustand, but that felt heavyweight. Instead, I:

\begin{itemize}
    \item Cache DOM elements at startup
    \item Use scoped variables (\texttt{isLoading}, \texttt{distributionChart}) within the event listener closure
    \item Keep async operations isolated with proper try/catch/finally blocks
    \item Render incrementally (don't rewrite the whole DOM)
\end{itemize}

It's not a framework, but it's \textbf{predictable and debuggable}.

\subsection{API Response Harmony}

Notice how every endpoint returns a \texttt{\{``status'': ``success''\}} blob? This is intentional. On the frontend, I can:

\begin{itemize}
    \item Always check for \texttt{data.status}
    \item Catch both HTTP errors and application-level errors
    \item Provide consistent error messaging to users
\end{itemize}

Small touch, huge impact on robustness.

\section{Common Pitfalls and How I Dodged Them}

\subsection{Pitfall 1: Hardcoding Dataset Paths}

I \textbf{almost} just put:

\begin{minted}{python}
df = pd.read_csv(r'C:\Users\Fofan\Downloads\archive\spam.csv')
\end{minted}

That works until you clone the repo on a different machine or reorganize your Downloads folder. Instead, I implemented path-checking logic that looks in multiple places. Small investment, prevents major headaches.

\subsection{Pitfall 2: Blocking Model Training}

Initially, the model trained synchronously on every Flask request. Slow clients would disconnect, thinking the server crashed. Now, it trains once on startup with proper error handling.

\begin{techdetail}
    In production, this would be a Celery task or similar background worker. Don't train ML models on your request thread!
\end{techdetail}

\subsection{Pitfall 3: Losing User Input on Error}

Early versions would clear the textarea when an error occurred. That's awful UX---users have to re-type everything. Now I preserve the input and just show an error message.

\section{Metrics That Actually Matter}

Let me be transparent about what this model achieves:

\begin{highlight}
    \textbf{Accuracy}: 98.39\% \\
    \textbf{Dataset}: 5,572 messages (4,825 ham, 747 spam) \\
    \textbf{Average ham length}: 71 characters \\
    \textbf{Average spam length}: 138.9 characters \\
    \textbf{Feature Extraction}: Bag-of-words with CountVectorizer \\
    \textbf{Classifier}: Multinomial Naive Bayes
\end{highlight}

\begin{devnote}
    98.39\% sounds incredible, right? But consider: if your model just predicted ``ham'' for everything, it'd be 86.6\% accurate (since ham is the majority class). So we're really only 11.7 percentage points better than a baseline. That's actually meaningful, but not revolutionary. Still, for binary classification on this data, Naive Bayes performs well.
\end{devnote}

\section{Visualizations and Dashboard Results}

\subsection{Dataset Distribution Visualization}

The doughnut chart displays the class imbalance in our training data:

\begin{itemize}
    \item \textbf{Safe Messages (Ham)}: 4,825 messages (86.6\%)
    \item \textbf{Spam Messages}: 747 messages (13.4\%)
\end{itemize}

This imbalance is typical in real-world datasets. We handle it by relying on accuracy metrics that account for the skew, and the model's strong performance suggests it learned to identify spam patterns effectively despite the minority class.

\begin{highlight}
    \textbf{Visualization Insight}: The doughnut chart is rendered client-side using Chart.js, allowing real-time updates without page refreshes. This keeps the UI feeling responsive and modern.
\end{highlight}

\subsection{Model Performance Dashboard}

The frontend displays four key metrics:

\subsubsection{Accuracy Card}

\begin{center}
\begin{tcolorbox}[colback=white, colframe=accentpurple, title=\textbf{Model Accuracy}, width=8cm]
    \centering\Large\textbf{98.39\%}
    
    \small\textit{On test set (1,115 holdout messages)}
\end{tcolorbox}
\end{center}

\subsubsection{Total Messages Card}

\begin{center}
\begin{tcolorbox}[colback=white, colframe=success, title=\textbf{Total Training Messages}, width=8cm]
    \centering\Large\textbf{5,572}
    
    \small\textit{Dataset size in production}
\end{tcolorbox}
\end{center}

\subsubsection{Message Length Analytics}

\textbf{Safe Messages (HAM)}:
\begin{center}
\begin{tcolorbox}[colback=white, colframe=success, width=8cm, height=2cm]
    \Large\textbf{Average: 71.0 characters}
    
    \small Legitimate messages tend to be concise
\end{tcolorbox}
\end{center}

\textbf{Spam Messages}:
\begin{center}
\begin{tcolorbox}[colback=white, colframe=danger, width=8cm, height=2cm]
    \Large\textbf{Average: 138.9 characters}
    
    \small Spam typically contains lengthy calls-to-action
\end{tcolorbox}
\end{center}

\subsection{User Interface Features}

The application provides a modern, glassmorphism-style interface built with pure CSS and vanilla JavaScript:

\begin{itemize}
    \item \textbf{Dark Mode Theme}: Reduces eye strain and looks professional
    \item \textbf{Real-time Results}: Displays prediction and confidence immediately upon submission
    \item \textbf{Gradient Backgrounds}: Uses CSS gradients and backdrop filters for depth
    \item \textbf{Interactive Charts}: Chart.js doughnut chart for dataset distribution
    \item \textbf{Confidence Visualization}: Visual progress bar showing model confidence level
\end{itemize}

\begin{devnote}
    The UI design respects dark mode conventions: using muted purples, greens, and reds against a dark background. This color scheme is WCAG AA compliant and looks clean without being sterile.
\end{devnote}

\subsection{Real-time Prediction Results}

Here's what users see when they submit a message:

\paragraph{Safe Message Result:}
\begin{center}
\begin{tcolorbox}[colback=white, colframe=success, title=\textbf{[OK] Safe Message}]
    This message appears to be legitimate and secure. No spam indicators detected.
    
    \textbf{Detection Confidence}: 94\%
    
    \textcolor{success}{\rule{0.75\textwidth}{5pt}}
\end{tcolorbox}
\end{center}

\paragraph{Spam Message Result:}
\begin{center}
\begin{tcolorbox}[colback=white, colframe=danger, title=\textbf{[WARN] Spam Detected}]
    This message has been flagged as potential spam. Exercise caution with links or requests.
    
    \textbf{Detection Confidence}: 89\%
    
    \textcolor{danger}{\rule{0.75\textwidth}{5pt}}
\end{tcolorbox}
\end{center}

\subsection{API Response Format}

Behind the scenes, the frontend receives structured JSON responses:

\begin{minted}{json}
{
  "prediction": "SPAM",
  "confidence": 92.45,
  "status": "success"
}
\end{minted}

This clean format allows the frontend to:
\begin{itemize}
    \item Display the prediction label (SPAM or HAM)
    \item Render a visual confidence bar
    \item Handle errors gracefully if status is not ``success''
\end{itemize}

\subsection{Dashboard Layout}

The dashboard is organized into logical sections:

\begin{enumerate}
    \item \textbf{Input Section}: Large textarea for message submission with character counter
    \item \textbf{Result Section}: Hidden until prediction is complete, shows instantly with animation
    \item \textbf{Stats Section}: Model accuracy and total messages processed
    \item \textbf{Chart Section}: Doughnut chart showing dataset distribution
    \item \textbf{Insights Section}: Average message lengths for spam vs. ham
\end{enumerate}

All sections are animated on page load for a polished, professional feel.

\section{What I'd Do Differently}

If I were starting over:

\begin{enumerate}
    \item \textbf{Add logging infrastructure early}. I'm using print statements, which work fine locally but become chaos in production.
    \item \textbf{Use environment variables} for configuration instead of hardcoding paths.
    \item \textbf{Build an admin dashboard} to manually test predictions and collect feedback.
    \item \textbf{Implement SQLite tracking} to log all predictions for later model retraining.
    \item \textbf{Use pytest} from day one. Testing this by hand gets tedious fast.
\end{enumerate}

\section{Wrapping Up}

This project taught me that sometimes the best stack isn't the trendiest one. Flask + Vanilla JS + Scikit-Learn is \textbf{boring in the best way}---it just works, it's maintainable, and it gets the job done.

The real wins here weren't fancy algorithms or bleeding-edge frameworks. They were:

\begin{itemize}
    \item Graceful error handling
    \item Real, meaningful metrics on the dashboard
    \item Thoughtful API design
    \item Simple, understandable frontend code
\end{itemize}

If you're building an MVP, sometimes the simplest choice is the right one. Ship it, gather feedback, and iterate.

\begin{thebibliography}{99}
\bibitem{sklearn} Pedregosa et al., \textit{Scikit-learn: Machine Learning in Python}, JMLR, 2011.
\bibitem{flask} Ronacher, \textit{Flask: A Modern Web Framework}, 2010--2024.
\bibitem{chartjs} Evrard, \textit{Chart.js}, JavaScript Charting Library, 2013--2024.
\end{thebibliography}

\end{document}
